\documentclass[10pt]{article}

\usepackage[margin=0.7in]{geometry}
\usepackage{amsmath,amssymb}
\usepackage{microtype}

\newcommand{\T}{\mathcal{T}}
\newcommand{\C}{\mathbb{C}}
\newcommand{\Rnn}{\mathbb{R}_{\ge 0}}
\newcommand{\Hh}{\mathcal{H}}
\newcommand{\Hv}{\mathcal{H}_V}

\setlength{\parindent}{1em}
\setlength{\parskip}{0.18em}
\setlength{\abovedisplayskip}{5pt}
\setlength{\belowdisplayskip}{5pt}

\title{\vspace{-1.6em}Adaptive Catamorphisms for Weighted Operational Theories}
\author{}
\date{\vspace{-2em}}

\begin{document}
\maketitle
\vspace{-1em}

\paragraph{Adaptive weighting.}
Assume the weighted operational-theory construction developed previously.
Let $\Gamma R=\T(1,R)$ be the set of closed rewrites, let $U$ be a rig of
weights, and let $V$ be a presentation language for admissible weighting
catamorphisms.  An element $v\in V$ determines the current interpretation of
rewrite constructors.  To update that interpretation after a rewrite, introduce
\[
\delta:V\times \Gamma R\longrightarrow U\times V,
\qquad
\delta(v,r)=\bigl(w(v,r),\tau(v,r)\bigr).
\]
Here $w(v,r)$ is the weight assigned to $r$ by the current catamorphism and
$\tau(v,r)$ is the presentation to be used after $r$.  The catamorphism itself
need not be stateful while traversing the syntax of $r$: one first evaluates
$r$ using $v$, then updates $v$ between small-step reductions.

There are two distinct ways to incorporate $V$ into the dynamics.  In an
\emph{internally adaptive} theory, the presentation is part of the
configuration.  A rewrite
\[
r:p\rightsquigarrow q
\]
induces
\[
(p,v)\rightsquigarrow \bigl(q,\tau(v,r)\bigr)
\]
with weight $w(v,r)$.  In an \emph{externally adaptive} theory, by contrast,
the system evolves under a single presentation shared by all branches during
a step.  The current $v_n$ determines a weighted rewrite operator $A_{v_n}$,
and a separate classical control process replaces $v_n$ by $v_{n+1}$ between
steps.  The distinction is mild for stochastic semantics but crucial for
quantum interference.

\section*{Nonnegative weights}

Take $U=\Rnn$.  For a fixed presentation $v$, the weighted rewrite operator is
determined by
\[
A_v(p,q)
=
\sum_{\substack{r:s(r)=p\\t(r)=q}} w(v,r).
\]
Assuming local summability,
\[
\sum_{r:s(r)=p}w(v,r)<\infty,
\]
one may normalize the outgoing weights to obtain transition probabilities.

Internal adaptation gives a Markov process on the enlarged state space
$\Gamma P\times V$.  From $(p,v)$, the rewrite $r$ occurs with probability
proportional to $w(v,r)$ and produces the new state
$\bigl(t(r),\tau(v,r)\bigr)$.  Thus the future transition law may depend on the
rewrite history through the evolving presentation.  This is natural when the
weighting policy is intended to learn, accumulate resources, change priorities,
or otherwise respond to individual reductions.

External adaptation instead uses a sequence $v_0,v_1,v_2,\ldots$ chosen
independently of which stochastic branch is taken during a particular step.
The process distribution evolves as
\[
\mu_{n+1}=K_{v_n}\mu_n,
\]
where $K_{v_n}$ is the normalized transition kernel determined by $v_n$.
This is a time-inhomogeneous Markov process.  External adaptation is appropriate
when a scheduler, optimization loop, or surrounding block changes the weighting
policy globally between rounds rather than as a consequence of one particular
rewrite.

\section*{Complex amplitudes}

Now take $U=\C$.  For each presentation $v$, let $\Hh$ be the Hilbert space
with basis $\{|p\rangle:p\in\Gamma P\}$ and define
\[
A_v|p\rangle
=
\sum_{r:s(r)=p}w(v,r)\,|t(r)\rangle.
\]
Parallel rewrites with the same source and target contribute to the same basis
vector and therefore interfere.  Assume local summability and boundedness.
Choose $\alpha_v\geq\|A_v\|$ and put $B_v=A_v/\alpha_v$, so that $B_v$ is a
contraction.  Let $U_v$ be a unitary dilation of $B_v$ on
$\Hh\otimes\C^2$ satisfying
\[
U_v\bigl(|\psi\rangle\otimes|0\rangle\bigr)
=
B_v|\psi\rangle\otimes|0\rangle
+|\eta_{v,\psi}\rangle\otimes|1\rangle
\]
for some $|\eta_{v,\psi}\rangle$.  The ancillary qubit records whether the
dilated evolution remains in the block representing the original weighted
rewrite operator.

In the externally adaptive quantum theory, $v_n$ remains classical and is
shared by every coherent branch.  At step $n$, adjoin a fresh ancillary qubit
$|0\rangle_n$ and apply $U_{v_n}$ to the process register and that ancilla.
After $n$ steps, projection onto the all-zero ancillary register gives the
unnormalized branch
\[
\frac{A_{v_{n-1}}\cdots A_{v_1}A_{v_0}}
{\alpha_{v_{n-1}}\cdots\alpha_{v_1}\alpha_{v_0}}
|\psi_0\rangle\otimes|0^n\rangle.
\]
Equivalently, one may postselect the fresh ancilla after every step.  The
presentation is updated externally between applications of the dilations, so
it carries no which-path information.  Consequently all rewrites arriving at
the same process basis vector still interfere.  If each $A_{v_n}$ is already
unitary, one may take $\alpha_{v_n}=1$ and dispense with the dilation ancillas.
Because $v_n$ is external, its update cannot depend on which coherent rewrite
``occurred'' without introducing a measurement or other branch-distinguishing
interaction.

Internal adaptation is different.  Let $\Hv$ have basis states $|v\rangle$
for semantic presentation states, and define the branch-sensitive operator
$\widetilde A$ on $\Hh\otimes\Hv$ by
\[
\widetilde A|p,v\rangle
=
\sum_{r:s(r)=p}
w(v,r)\,|t(r),\tau(v,r)\rangle.
\]
This operator too must be bounded.  After rescaling to a contraction it is
implemented by a unitary dilation on
$\Hh\otimes\Hv\otimes\C^2$, again using a fresh dilation ancilla at each
small-step reduction.  The $V$ register and the dilation ancilla play distinct
roles: the ancilla implements a possibly nonunitary weighted operator, whereas
$V$ stores branch-dependent future semantics.

Suppose two rewrites
\[
r_1,r_2:p\rightsquigarrow q
\]
have amplitudes $a_1,a_2$ and update the presentation to $v_1,v_2$.  On the
postselected dilation branch their contribution is
\[
a_1|q,v_1\rangle+a_2|q,v_2\rangle.
\]
If $v_1=v_2$, the amplitudes combine to $(a_1+a_2)|q,v_1\rangle$.  If $v_1$
and $v_2$ are orthogonal presentation states, the $V$ register retains
which-path information and the histories do not interfere when only the
process register is observed.  More generally, the interference term is
scaled by $\langle v_2|v_1\rangle$.  Interference can return if the presentation
information is later coherently erased or uncomputed so that the branches
again occupy the same total state.

Because $V$ is a language of presentations, semantically equivalent
presentations should not be made distinguishable merely because their syntax
differs.  One may quotient $V$ by equivalence of the catamorphisms they denote,
or choose canonical representatives.  What matters for interference is
whether the retained adaptive state encodes distinguishable future behavior.

\paragraph{Summary.}
For $U=\Rnn$, external adaptation gives a globally time-varying transition
kernel, while internal adaptation gives a Markov process on processes paired
with adaptive state.  For $U=\C$, both constructions are implemented through
unitary dilations when the weighted operators are not themselves unitary.
External adaptation keeps $V$ outside the quantum state and therefore preserves
ordinary interference across rewrite paths.  Internal adaptation makes $V$ a
quantum register; branch-dependent updates then act as which-path information
unless that information is subsequently uncomputed.

\end{document}
